% GENERATED FILE. Edit canonical lessons, then run npm run generate:books.
% canonical-source-hash: fnv1a64:3c077a576f4e0922
% canonical-lessons: SA-C69-eye, SA-C69-throat, SA-C69-belly, SA-C69-forehead, SA-C69-tongue

\chapter{Eye, Throat and Tongue}
\label{ch:sa-eye-throat-tongue}
\hlchaptermodality{69}

\begin{chapteropening}
\textbf{By the end of this chapter:} \emph{I can name the eye, the throat, the belly, the forehead and the tongue.}

Complete the last lesson of chapter 69: i can name the eye, the throat, the belly, the forehead and the tongue.
\end{chapteropening}

\section[netram]{\sk{नेत्रम्} (netram) — an eye}

\label{lesson:SA-C69-eye}



\begin{quote}
Before the new one: say the Sanskrit for shyness, then the Sanskrit for joy.
\end{quote}

\subsection*{You'll want to know: \sk{नेत्रम्}}
\textbf{\sk{नेत्रम्}} (\emph{netram}) — \textquotedblleft{}an eye\textquotedblright{}, the word of poems, beside the everyday \textbf{\sk{अक्षि}}.

\begin{grammarlens}[title={What you've built}]
A second word for the eye.
\end{grammarlens}

\subsection*{Your turn}
\begin{itemize}
\raggedright
  \item \emph{Say it:} \emph{netram}
  \item \emph{Say it:} \emph{netram}, once more
  \item \emph{Say it:} say \emph{netram}
  \item \emph{Recall:} say the Sanskrit for shyness, then the Sanskrit for joy, then say \emph{netram} again
\end{itemize}

\begin{tcolorbox}[breakable,colback=teal!4,colframe=teal!35!black,title={Before you move on}]
What does \sk{नेत्रम्} mean? (\textquotedblleft{}An eye\textquotedblright{}.) Say it once more.
\end{tcolorbox}
\section[kaṇṭhaḥ]{\sk{कण्ठः} (kaṇṭhaḥ) — the throat}

\label{lesson:SA-C69-throat}



\begin{quote}
Before the new one: say the Sanskrit for joy, then the Sanskrit for an eye.
\end{quote}

\subsection*{You'll want to know: \sk{कण्ठः}}
\textbf{\sk{कण्ठः}} (\emph{kaṇṭhaḥ}) — \textquotedblleft{}the throat\textquotedblright{}. With a \textbf{\sk{कासः}}, it is the \textbf{\sk{कण्ठः}} that hurts.

\begin{grammarlens}[title={What you've built}]
Where the voice comes from.
\end{grammarlens}

\subsection*{Your turn}
\begin{itemize}
\raggedright
  \item \emph{Say it:} \emph{kaṇṭhaḥ}
  \item \emph{Say it:} \emph{kaṇṭhaḥ}, once more
  \item \emph{Say it:} say \emph{kaṇṭhaḥ}
  \item \emph{Recall:} say the Sanskrit for joy, then the Sanskrit for an eye, then say \emph{kaṇṭhaḥ} again
\end{itemize}

\begin{tcolorbox}[breakable,colback=teal!4,colframe=teal!35!black,title={Before you move on}]
What does \sk{कण्ठः} mean? (\textquotedblleft{}The throat\textquotedblright{}.) Say it once more.
\end{tcolorbox}
\section[udaram]{\sk{उदरम्} (udaram) — the belly}

\label{lesson:SA-C69-belly}



\begin{quote}
Before the new one: say the Sanskrit for an eye, then the Sanskrit for the throat.
\end{quote}

\subsection*{You'll want to know: \sk{उदरम्}}
\textbf{\sk{उदरम्}} (\emph{udaram}) — \textquotedblleft{}the belly\textquotedblright{}, where \textbf{\sk{क्षुधा}} is felt.

\begin{grammarlens}[title={What you've built}]
Where hunger lives.
\end{grammarlens}

\subsection*{Your turn}
\begin{itemize}
\raggedright
  \item \emph{Say it:} \emph{udaram}
  \item \emph{Say it:} \emph{udaram}, once more
  \item \emph{Say it:} say \emph{udaram}
  \item \emph{Recall:} say the Sanskrit for an eye, then the Sanskrit for the throat, then say \emph{udaram} again
\end{itemize}

\begin{tcolorbox}[breakable,colback=teal!4,colframe=teal!35!black,title={Before you move on}]
What does \sk{उदरम्} mean? (\textquotedblleft{}The belly\textquotedblright{}.) Say it once more.
\end{tcolorbox}
\section[lalāṭam]{\sk{ललाटम्} (lalāṭam) — the forehead}

\label{lesson:SA-C69-forehead}



\begin{quote}
Before the new one: say the Sanskrit for the throat, then the Sanskrit for the belly.
\end{quote}

\subsection*{You'll want to know: \sk{ललाटम्}}
\textbf{\sk{ललाटम्}} (\emph{lalāṭam}) — \textquotedblleft{}the forehead\textquotedblright{}, above the \textbf{\sk{अक्षि}}.

\begin{grammarlens}[title={What you've built}]
The top of the face.
\end{grammarlens}

\subsection*{Your turn}
\begin{itemize}
\raggedright
  \item \emph{Say it:} \emph{lalāṭam}
  \item \emph{Say it:} \emph{lalāṭam}, once more
  \item \emph{Say it:} say \emph{lalāṭam}
  \item \emph{Recall:} say the Sanskrit for the throat, then the Sanskrit for the belly, then say \emph{lalāṭam} again
\end{itemize}

\begin{tcolorbox}[breakable,colback=teal!4,colframe=teal!35!black,title={Before you move on}]
What does \sk{ललाटम्} mean? (\textquotedblleft{}The forehead\textquotedblright{}.) Say it once more.
\end{tcolorbox}
\section[jihvā]{\sk{जिह्वा} (jihvā) — the tongue}

\label{lesson:SA-C69-tongue}



\begin{quote}
Before the new one: say the Sanskrit for the belly, then the Sanskrit for the forehead.
\end{quote}

\subsection*{You'll want to know: \sk{जिह्वा}}
\textbf{\sk{जिह्वा}} (\emph{jihvā}) — \textquotedblleft{}the tongue\textquotedblright{}, inside the \textbf{\sk{मुखम्}}.

\begin{grammarlens}[title={What you've built}]
What tastes and speaks.
\end{grammarlens}

\subsection*{Your turn}
\begin{itemize}
\raggedright
  \item \emph{Say it:} \emph{jihvā}
  \item \emph{Say it:} \emph{jihvā}, once more
  \item \emph{Say it:} say \emph{jihvā}
  \item \emph{Recall:} say the Sanskrit for the belly, then the Sanskrit for the forehead, then say \emph{jihvā} again
\end{itemize}

\begin{tcolorbox}[breakable,colback=teal!4,colframe=teal!35!black,title={Before you move on}]
What does \sk{जिह्वा} mean? (\textquotedblleft{}The tongue\textquotedblright{}.) Say it once more.
\end{tcolorbox}
